\documentclass[11pt,a4paper]{article}
\usepackage[T1]{fontenc}
\usepackage[utf8]{inputenc}
\usepackage{amsmath,amssymb,geometry,hyperref}
\geometry{margin=25mm}
\setlength{\emergencystretch}{2em}
\hypersetup{colorlinks=true,urlcolor=blue,linkcolor=blue}
\title{Universe Tunnel Lab\\\large A playful picture, with inspectable mathematics}
\author{Concept: Victor}
\date{31 October 2026}
\begin{document}
\maketitle
\begin{abstract}
Imagine the history of a universe as a tunnel: each slice stands for three-dimensional space, and extra visual features carry information about matter, fields and correlations. This educational project explores that picture using small, explicitly specified models. A tunnel is a visualization convention; physical predictions come from the equations and assumptions below.
\end{abstract}
\section{Three directions, many state variables}
A spatial dimension describes an independent direction in physical space. A state-space dimension describes an independent parameter needed to specify a state. They are different mathematical roles. A two-qubit pure state has six real parameters after normalization and removal of global phase; a general mixed state has fifteen. These are not extra directions in a room.

For an appropriate classical-geometry/quantum-matter approximation, a schematic slice record is
\begin{equation}
\mathcal U_t=(\Sigma_t,h_{ij},K_{ij},\rho_t),\qquad
\mathcal T=\{\mathcal U_t:t\in I\}.
\end{equation}
Here $\Sigma_t$ is a spatial hypersurface, $h_{ij}$ its metric, $K_{ij}$ its extrinsic curvature, and $\rho_t$ the density operator of retained matter/field degrees of freedom. A classical field model instead stores the necessary field values and conjugate momenta. Constraints restrict the allowable data. A global foliation requires appropriate spacetime assumptions, such as global hyperbolicity. This tuple alone supplies neither evolution laws nor quantum gravity.

The tube wall displays selected observables. Interior links display correlations. Neither part is a separate physical domain; quantum mechanics describes both. Tunnel length is a chosen time label. There is no second external time moving the universe along the diagram.
\section{The two-qubit laboratory}
Two fixed local sites are at $\mathbf r_A=(-d/2,0,0)$ and $\mathbf r_B=(d/2,0,0)$. Exactly one excitation is retained, in the ordered basis $\{|10\rangle,|01\rangle\}$. An effective exchange Hamiltonian is
\begin{equation}
H=\frac{\hbar J}{2}(X_AX_B+Y_AY_B).
\end{equation}
The symbols $X,Y,Z$ denote Pauli operators, not spatial coordinates. Starting from $|10\rangle$, isolated proper-time evolution gives
\begin{equation}
|\psi(\tau)\rangle=\cos(J\tau)|10\rangle-i\sin(J\tau)|01\rangle.
\end{equation}
Consequently $P_A=\cos^2(J\tau)$, $P_B=\sin^2(J\tau)$ and concurrence $C=|\sin(2J\tau)|$.

The implemented dephasing model is
\begin{equation}
\frac{d\rho}{d\tau}=-\frac{i}{\hbar}[H,\rho]
+\frac{\gamma}{4}\sum_{k=A,B}(Z_k\rho Z_k-\rho).
\end{equation}
This is a reduced, Markovian open-system model. It describes phase noise, not energy relaxation or a complete thermal environment. Its convention makes the $|10\rangle\langle01|$ coherence decay at rate $\gamma$ in the absence of exchange.

In the one-excitation subspace, write
\begin{equation}
\rho=\frac12\begin{pmatrix}1+z&x-iy\\x+iy&1-z\end{pmatrix}.
\end{equation}
The Bloch variables $x,y,z$ are state coordinates, not positions. Direct substitution gives
\begin{equation}
\dot x=-\gamma x,\qquad\dot y=-\gamma y-2Jz,\qquad\dot z=2Jy,
\end{equation}
with dots denoting proper-time derivatives. Thus
\begin{equation}
\ddot z+\gamma\dot z+4J^2z=0.
\end{equation}
The engine uses analytic damped-oscillator solutions, including critical and overdamped limits. In this restricted state family,
\begin{equation}
P_A=(1+z)/2,\quad P_B=(1-z)/2,\quad
C=\sqrt{x^2+y^2},\quad \operatorname{Tr}\rho^2=(1+x^2+y^2+z^2)/2.
\end{equation}
These concurrence and purity expressions should not be reused for arbitrary two-qubit states.

Three initial states are provided: $(x,y,z)=(0,0,1)$ for $|10\rangle$; $(1,0,0)$ for $(|10\rangle+|01\rangle)/\sqrt2$; and $(0,0,0)$ for an equal mixture of the two basis states. The last two have identical local probabilities but different entanglement. All retain $\langle Z_AZ_B\rangle=-1$: no entanglement does not mean no correlation.
\section{Clock comparisons and simultaneity}
The local dynamics use $\tau=\alpha t$ with a constant clock ratio. The three selectable models are
\begin{equation}
\alpha=1,\qquad
\alpha=\sqrt{1-v^2/c^2},\qquad
\alpha=\sqrt{1-2GM/(rc^2)}.
\end{equation}
The SR comparison is a small inertial lab. The GR comparison is a small lab held at fixed exterior Schwarzschild radius, relative to time at infinity, outside both the material source and horizon. Tidal variation and backreaction are neglected. The settings are alternatives; their factors are not combined into a claimed general solution. The effective exchange model does not resolve the propagation time of a mediator.

The separate relativity diagram uses $T=ct/L$, $X=x/L$, $\beta=v/c$ and $\Gamma=(1-\beta^2)^{-1/2}$:
\begin{align}
T'&=\Gamma(T-\beta X),&X'&=\Gamma(X-\beta T),\\
T&=\beta X+T'/\Gamma,&-T'^2+X'^2&=-T^2+X^2.
\end{align}
A line of constant $T'$ is a different simultaneity slice for each velocity. At $\beta=0.6$, the events $(T,X)=(0,-1)$ and $(0,+1)$ have $T'=+0.75$ and $-0.75$. They are spacelike separated. Their coordinate ordering can differ without reversing a causal influence.
\section{An evolving cosmological background}
A homogeneous, isotropic FLRW model has
\begin{equation}
ds^2=-c^2dt^2+a^2(t)d\ell_k^2.
\end{equation}
The reference spatial metric $d\ell_k^2$ has length units; $a$ is dimensionless. Cosmic time is proper time of comoving observers. It is useful in this symmetry reduction, not a universal simultaneity imposed on every observer.

The reference-density parameters obey $\Omega_k=1-\Omega_m-\Omega_r-\Omega_\Lambda$. With $H=\dot a/a$,
\begin{equation}
\frac{H^2}{H_{\rm ref}^2}=\frac{\Omega_r}{a^4}+\frac{\Omega_m}{a^3}+\frac{\Omega_k}{a^2}+\Omega_\Lambda.
\end{equation}
Set $u=H_{\rm ref}(t-t_{\rm start})$ and $v=da/du$. The numerical system is
\begin{equation}
a'=v,\qquad v'=-\frac{\Omega_r}{a^3}-\frac{\Omega_m}{2a^2}+\Omega_\Lambda a,
\qquad\chi'=\frac1a.
\end{equation}
Its initial positive $v$ is fixed by the Friedmann constraint. Evolving acceleration allows $v$ to pass through zero into contraction. Adaptive RK4 with step doubling and Richardson correction is checked against analytic solutions and independent high-accuracy integration. The default display spans $a_{\rm start}=0.2$ to $a=3$, $a=0.05$ during contraction, or $u=8$, whichever occurs first. These are computational limits.

Matter and radiation are separately conserved ideal fluids, with density scaling as $a^{-3}$ and $a^{-4}$; cosmological-constant density is constant. Matter combines baryons and cold dark matter into one pressureless component. Curvature affects geometry; it is excluded from the material-density composition bar.

A test photon emitted at the first slice has
\begin{equation}
\frac{\lambda}{\lambda_{\rm emit}}=\frac{a}{a_{\rm start}},\qquad
\frac{E}{E_{\rm emit}}=\frac{a_{\rm start}}{a},\qquad E=\frac{hc}{\lambda}.
\end{equation}
Multiply $\chi=\int du/a$ by $c/H_{\rm ref}$ to obtain accumulated comoving path length. It is not the full particle horizon or, in a closed geometry, the shortest distance to the emission point. The drawn sinusoid illustrates wavelength; it is not a Maxwell field solution. A bound local apparatus is not automatically stretched by $a(t)$.
\section{Where more physics would go}
The Standard Model is a quantum field theory with gauge structure
\begin{equation}SU(3)_c\times SU(2)_L\times U(1)_Y.\end{equation}
The strong interaction is associated with $SU(3)_c$. Electromagnetism is the unbroken subgroup after electroweak symmetry breaking. Gravity is described by GR, and quantum mechanics is a framework used by the SM, so these labels are not independent forces to add as one-number axes.

An established low-energy effective-field-theory organization, in units $\hbar=c=1$, is schematically
\begin{equation}
S_{\rm eff}=\int d^4x\sqrt{-g}\left[
\frac{M_{\rm Pl}^2}{2}(R-2\Lambda)+\mathcal L_{\rm SM}
+\alpha R^2+\beta_R R_{\mu\nu}R^{\mu\nu}+\cdots\right].
\end{equation}
$M_{\rm Pl}$ is the reduced Planck mass. The dots include further symmetry-allowed curvature and matter operators. This is an organizing framework, not a new equation of this project or a complete high-energy quantum-gravity theory. The demos do not solve this action. Supersymmetry is an optional theoretical extension, not implied by rhetorical opposites.

A sensible next module is linear density growth on the cosmological background, followed by a specified electromagnetic field model. Each extension needs units, constraints, evolution laws and known limiting cases. A finite quantum laboratory, homogeneous cosmos and flat-spacetime coordinate diagram are currently separate experiments; their shared visual language is not a claim of full dynamical unification.
\section*{References and provenance}
Concept and visual starting point: Victor. Implementation and explanatory draft developed with AI assistance; corrections and independent checks are welcome. Existing mathematical precedents include state spaces, spacetime foliations and correlation visualizations. No novelty or priority claim is made for the tunnel metaphor.
\begin{enumerate}
\item E. Gourgoulhon, \emph{3+1 Formalism and Bases of Numerical Relativity}, \url{https://arxiv.org/abs/gr-qc/0703035}.
\item Einstein Online, \emph{Spacetime}, \url{https://www.einstein-online.info/en/spacetime/}.
\item W. K. Wootters, \emph{Entanglement of Formation of an Arbitrary State of Two Qubits}, \url{https://arxiv.org/abs/quant-ph/9709029}.
\item W. H. Zurek, \emph{Decoherence, einselection, and the quantum origins of the classical}, \url{https://arxiv.org/abs/quant-ph/0105127}.
\item D. Tong, \emph{Cosmology: The Expanding Universe}, \href{https://davidtong.org/teaching/cosmology/cosmohtml/S1}{Course notes, section 1}.
\item CERN, \emph{The Standard Model}, \url{https://home.cern/science/physics/standard-model}.
\item J. F. Donoghue, \emph{Quantum General Relativity and Effective Field Theory}, \url{https://arxiv.org/abs/2211.09902}.
\end{enumerate}
\end{document}
